\chaptercustom{5}{TRÍCH DẪN, THAM CHIẾU VÀ LIÊN KẾT}
Chương này hướng dẫn trích dẫn tài liệu tham khảo, tham chiếu chéo tới hình, bảng, phương trình, tạo liên kết, và liệt kê các lỗi biên dịch thường gặp.

\subsection{Trích dẫn tài liệu tham khảo}
\label{sec:trichdan}
Tài liệu tham khảo được khai báo trong tệp \path{KetThuc/TaiLieuThamKhao.bib}. Mỗi tài liệu là một mục BibTeX, có thể lấy từ Google Scholar (nút \textbf{Trích dẫn} $\rightarrow$ \textbf{BibTeX}):
\begin{verbatim}
@book{howell2016principles,
  title={Principles of Fourier analysis},
  author={Howell, Kenneth B},
  year={2016},
  publisher={CRC Press}
}
\end{verbatim}
Trong bài, trích dẫn bằng \verb|\cite{khóa}|, ví dụ \verb|\cite{bracewell1989fourier}| cho ra \cite{bracewell1989fourier}. Trích dẫn nhiều tài liệu cùng lúc thì liệt kê các khóa, cách nhau bằng dấu phẩy, ví dụ \cite{stein2011fourier, howell2016principles}. Kiểu IEEE tự đánh số theo thứ tự trích dẫn, và chỉ những tài liệu đã được \verb|\cite| mới xuất hiện trong danh mục tài liệu tham khảo.

\subsection{Tham chiếu chéo}
Đặt nhãn bằng \verb|\label{tên}| ngay sau \verb|\caption| (với hình, bảng, thuật toán), bên trong phương trình, định lý, hoặc ngay sau đề mục; rồi dùng \verb|\ref{tên}| để in số. Ví dụ: Hình \ref{hinh:mach}, Bảng \ref{bang:thinghiem}, PT \ref{pt:dft}, Định lý \ref{dl:laymau}, Thuật toán \ref{tt:dft}, Mục \ref{sec:trichdan}. Số luôn đúng khi thêm hoặc bớt nội dung. Nên đặt tên nhãn có tiền tố theo loại (\texttt{hinh:}, \texttt{bang:}, \texttt{pt:}) để dễ nhớ.

\subsection{Liên kết}
Mục lục, danh mục, trích dẫn và tham chiếu đều bấm được trong tệp PDF. Liên kết ra ngoài viết bằng \verb|\href{địa chỉ}{chữ hiển thị}|, ví dụ \href{https://hust.edu.vn}{trang chủ HUST}. Mặc định liên kết không có màu để bản in gọn; muốn hiện khung màu khi soạn, trong \texttt{thesis-config.cls} đổi hai dòng sau thành:
\begin{verbatim}
% \RequirePackage[unicode,hidelinks]{hyperref}
\RequirePackage[unicode]{hyperref}
\end{verbatim}

\subsection{Lỗi thường gặp}
\begin{itemize}
    \item Số hiện thành \texttt{??}: biên dịch lại, hoặc kiểm tra tên nhãn trong \verb|\ref| có khớp với \verb|\label| không.
    \item \texttt{Undefined control sequence}: gõ sai tên lệnh, hoặc lệnh cần gói chưa được nạp.
    \item \texttt{File ... not found}: đường dẫn ảnh tính từ thư mục chứa \texttt{main.tex}; Overleaf phân biệt chữ hoa và chữ thường trong tên tệp.
    \item Các ký tự \verb|% & _ # $| có nghĩa đặc biệt trong LaTeX; muốn in ra thì viết \verb|\% \& \_ \# \$|.
\end{itemize}

\cleardoublepage
